% ============================================================================
% PATENTE DE INVENÇÃO — MANIPULAÇÃO DE CAMPO Ψ POR RESSONÂNCIA OH
% Instituto Nacional da Propriedade Industrial (INPI) — Brasil
% Natureza: Patente de Invenção (PI)
% Classificação IPC: G21B 3/00; H05H 1/22; G01V 7/00; H01L 39/00
% ============================================================================

\documentclass[12pt,a4paper]{article}

\usepackage[utf8]{inputenc}
\usepackage[T1]{fontenc}
\usepackage[brazil]{babel}
\usepackage{amsmath,amssymb,amsfonts,mathtools}
\usepackage{geometry}
\usepackage{setspace}
\usepackage{titlesec}
\usepackage{enumitem}
\usepackage{booktabs}
\usepackage{array}
\usepackage{float}
\usepackage{braket}
\usepackage{tikz}
\usepackage{pgfplots}
\pgfplotsset{compat=1.18}
\usetikzlibrary{shapes.geometric,arrows.meta,positioning,calc,
  decorations.pathreplacing,patterns,fit,backgrounds,
  decorations.markings,decorations.pathmorphing,shadings}

\geometry{left=3cm, right=2cm, top=3cm, bottom=2cm}
\onehalfspacing

\titleformat{\section}{\normalfont\Large\bfseries}{\thesection}{1em}{}
\titleformat{\subsection}{\normalfont\large\bfseries}{\thesubsection}{1em}{}
\titleformat{\subsubsection}{\normalfont\normalsize\bfseries}{\thesubsubsection}{1em}{}

% Comandos auxiliares
\newcommand{\alphaii}{\alpha^2}
\newcommand{\Dfolds}{D_{\text{folds}}}
\newcommand{\Fexp}{F_{\text{exp}}}
\newcommand{\thetaM}{\theta_{\text{Miguel}}}
\newcommand{\Hfour}{\mathcal{H}_4}
\newcommand{\rhoPsi}{\rho_\Psi}
\newcommand{\nuOH}{\nu_{\text{OH}}}
\newcommand{\tfix}{t_{\text{fix}}}
\newcommand{\Tcrit}{T_c}

\begin{document}

\begin{center}
{\Large\bfseries RELATÓRIO DESCRITIVO}\\[0.8cm]
{\large\bfseries MÉTODO E SISTEMA DE MANIPULAÇÃO DE CAMPO\\
PSIÔNICO POR RESSONÂNCIA ELETROMAGNÉTICA NA\\
FREQUÊNCIA DA HIDROXILA PARA INDUÇÃO DE\\
TRANSIÇÕES DE FASE EM MEIOS DENSOS}\\[0.8cm]
\textbf{Inventor:} Luiz Antonio Rotoli Miguel\\[0.3cm]
\textbf{Classificação IPC:} G21B 3/00; H05H 1/22;
  G01V 7/00; H01L 39/00\\[0.3cm]
\textbf{Referência a pedidos anteriores:}\\
BR 10 2025 026951 1 (ACOM --- Compressão Ontológica de Memória)\\
ACOM Mirror (Reflexão Dimensional de Dados)\\
Supercondutor Holográfico TGL (Supercondutividade por
  Ressonância Luminodinâmica)
\end{center}

\vspace{1cm}

% ============================================================================
\section{Campo da invenção}
% ============================================================================

A presente invenção refere-se a métodos e sistemas de manipulação
de campo psiônico ($\Psi$) por ressonância eletromagnética na
frequência do radical hidroxila (OH), aplicáveis aos setores de
energia nuclear, ciência de materiais, supercondutividade,
detecção astrofísica e física fundamental. Mais especificamente,
a invenção descreve um método para indução controlada de
transições de fase em meios densos (deuterados, metálicos,
supercondutores ou outros) mediante exposição simultânea a campo
eletromagnético sintonizado na frequência $\nuOH \approx 1{,}666$
GHz (frequência das linhas de emissão do radical hidroxila,
$\lambda \approx 18$ cm) e regime criogênico na temperatura
crítica $\Tcrit = h\nuOH / k_B \approx 80$ mK, provocando
nucleação de domínios de fase densa (``gotículas psiônicas'')
onde o acoplamento entre campo gravitacional e campo
eletromagnético, parametrizado pela Constante de Miguel
$\alphaii = 0{,}012031$, se torna operante, alterando
propriedades macroscópicas do meio.

\newpage
% ============================================================================
\section{Fundamentos da invenção}
% ============================================================================

\subsection{Estado da técnica}

\subsubsection{Fusão nuclear a baixas energias (LENR)}

Desde o anúncio de Fleischmann e Pons (1989) de excesso de calor
em células eletrolíticas de D$_2$O com cátodo de paládio,
numerosos grupos relataram fenômenos anômalos em sistemas
deuterados a baixa temperatura, porém sem reprodutibilidade
consistente e sem mecanismo teórico aceito. O estado da técnica
inclui: eletrólise em paládio deuterado (Fleischmann-Pons),
descarga luminescente em deutério (Karabut), bombardeamento
iônico de deutério em metais (Iwamura), e sonofusão
(Taleyarkhan). Nenhum desses métodos opera com mecanismo
físico estabelecido para redução da barreira Coulombiana.

\subsubsection{Confinamento magnético e inercial}

ITER (tokamak, confinamento magnético a $10^8$ K) e NIF
(confinamento inercial por laser) representam abordagens de
fusão ``quente'' com custos superiores a 10$^{10}$ USD e
temperaturas de plasma milhões de vezes acima da temperatura
ambiente. Não existem abordagens de fusão que operem em
regime criogênico ($<1$ K) com mecanismo de redução de
barreira via acoplamento gravitacional-eletromagnético.

\subsubsection{Transições de fase em meios densos}

O estudo de transições de fase em matéria condensada
(Landau-Ginzburg, BCS para supercondutividade, condensação de
Bose-Einstein) não incorpora acoplamento com campo gravitacional.
Nenhuma teoria existente prevê que irradiação em frequência
específica (1,666 GHz) acoplada a regime criogênico possa
induzir transições de fase com propriedades anômalas mediadas
por campo gravitacional.

\subsubsection{Detecção de matéria escura}

Experimentos de detecção direta de matéria escura (XENON1T,
LUX-ZEPLIN, PandaX) buscam recuo nuclear por interação com
WIMPs. Nenhum experimento existente utiliza deslocamento de
frequência de linhas espectrais como indicador de densidade
de campo escuro, nem emprega ressonância eletromagnética
na frequência da hidroxila como sonda de campo psiônico.

\subsubsection{Maser de hidroxila}

As linhas de emissão do radical OH em 1,6654018 GHz e
1,6673590 GHz (transições $\Lambda$-doubling) são conhecidas
em radioastronomia desde 1963 (Weinreb et al.). OH produz
masers naturais em nuvens moleculares, regiões de formação
estelar e envoltórios de estrelas evoluídas. As transições
$\Lambda$-doubling são extremamente sensíveis a campos
magnéticos, elétricos e, conforme a presente invenção, a
campos gravitacionais via acoplamento TGL.

\subsection{Problemas técnicos não resolvidos}

\textbf{Problema 1 --- Barreira Coulombiana.} Todos os métodos
conhecidos de fusão nuclear requerem temperaturas extremas
($>10^7$ K) ou confinamento mecânico para superar a repulsão
eletrostática entre núcleos. Não existe método que reduza a
barreira Coulombiana efetiva a baixas temperaturas.

\textbf{Problema 2 --- Ausência de mecanismo de acoplamento.}
Na física padrão, campo gravitacional e campo eletromagnético
são tratados como independentes (Relatividade Geral +
Eletrodinâmica Quântica). Não existe formulação que permita
campo eletromagnético em frequência específica alterar o
acoplamento gravitacional local.

\textbf{Problema 3 --- Manipulação de matéria escura.} O campo
psiônico ($\Psi$), identificado na TGL como substrato da
matéria escura, não possui mecanismo conhecido de manipulação
controlada em laboratório.

\textbf{Problema 4 --- Transições de fase sem frequência
de controle.} Transições de fase em matéria condensada são
tipicamente induzidas por variação de temperatura ou pressão,
sem frequência eletromagnética específica que sirva como
``canal de controle'' para o campo psiônico subjacente.

\subsection{Solução proposta}

A presente invenção resolve os quatro problemas ao identificar
a frequência do radical hidroxila ($\nuOH \approx 1{,}666$ GHz)
como a \textbf{frequência de ressonância do campo psiônico}
$\Psi$, e ao estabelecer um protocolo de três etapas
(ressonância, resfriamento, nucleação) que induz transições de
fase em meios densos com propriedades anômalas mediadas pelo
acoplamento gravitacional-eletromagnético da TGL.

O fundamento teórico reside na identificação:
\begin{equation}
\boxed{\omega_\Psi = 2\pi \times \nuOH = 2\pi \times 1{,}666
  \text{ GHz}}
\end{equation}
como frequência natural de oscilação do campo $\Psi$ (``água
escura''), derivada da correspondência entre o radical OH
(precursor molecular da água, H$_2$O = H + OH) e o campo
$\Psi$ (substrato psiônico em transição de fase, ``água
escura'' na terminologia TGL).

\newpage
% ============================================================================
\section{Objetivos da invenção}
% ============================================================================

\subsection{Objetivo principal}

Prover método e sistema de manipulação controlada do campo
psiônico $\Psi$ em meios densos, utilizando ressonância
eletromagnética na frequência da hidroxila ($\nuOH \approx
1{,}666$ GHz) acoplada a regime criogênico ($T \leq \Tcrit
\approx 80$ mK), para indução de transições de fase com
propriedades anômalas mediadas pelo acoplamento
gravitacional-eletromagnético parametrizado pela Constante
de Miguel $\alphaii = 0{,}012031$.

\subsection{Objetivos específicos}

(a)~Definir a ``Régua TGL'' --- conjunto de equações de escala
que determinam, para qualquer sistema material com frequência
característica $\nu^*$ e densidade $\rho$, a escala espacial
crítica $L^*$, a temperatura crítica $\Tcrit$, o tempo de
fixação gravitacional $\tfix$ e a condição de transição entre
regime químico e regime gravitacional.

(b)~Especificar o protocolo de três etapas (ressonância OH,
resfriamento criogênico, nucleação de gotículas) para indução
controlada de domínios de fase densa (``gotículas psiônicas'')
onde o acoplamento TGL é operante.

(c)~Descrever a sequência de colapso ontológico ``água escura
$\to$ OH $\to$ H$_2$O'' como manifestação física da
radicalização/ressurreição TGL, estabelecendo correspondência
entre o colapso do campo $\Psi$ em matéria ordinária e as
operações ACOM ($g = \sqrt{|L|}$, $s = \text{sign}(L)$,
$L' = s \times g^2$).

(d)~Prover sistema (aparato) capaz de: gerar campo EM
estável em $\nuOH$ com potência controlável; manter alvo
em temperatura $\leq 80$ mK; detectar produtos de reações
nucleares (nêutrons, gamas), anomalias térmicas e
deslocamentos espectrais; e monitorar em tempo real
indicadores de ativação do campo $\Psi$.

(e)~Especificar aplicações em pelo menos quatro domínios:
intensificação de reações nucleares em meios deuterados;
indução de estados supercondutores por coerência psiônica;
detecção de matéria escura por deslocamento de linhas OH;
e criação de materiais com propriedades anômalas por
nucleação de fase densa.

\newpage
% ============================================================================
\section{Descrição detalhada da invenção}
% ============================================================================

\subsection{Fundamento teórico --- A Régua TGL}

A Teoria da Gravitação Luminodinâmica (TGL) estabelece que o
campo gravitacional e o campo eletromagnético são manifestações
de um campo unificado, o campo psiônico $\Psi$, com constante
de acoplamento $\alphaii = 0{,}012031$ (Constante de Miguel).
A presente invenção aplica esse acoplamento a sistemas materiais
reais através de um conjunto de equações de escala denominado
``Régua TGL''.

\subsubsection{Frequência de ressonância do campo $\Psi$}

A frequência de ressonância do campo $\Psi$ é identificada com
a frequência do radical hidroxila (OH):
\begin{equation}
\nu_\Psi \equiv \nuOH \approx 1{,}666 \text{ GHz}
  \quad (\lambda \approx 18 \text{ cm})
\label{eq:freq_res}
\end{equation}

Esta identificação fundamenta-se em três correspondências:

\textbf{Correspondência molecular:} OH é o precursor da água
(H$_2$O = H + OH). Na TGL, o campo $\Psi$ é denominado ``água
escura'' por ser o substrato não-colapsado de que a matéria
ordinária (inclusive a água) é manifestação colapsada. A
frequência de OH é portanto a frequência do estado
pré-colapsado.

\textbf{Correspondência cosmológica:} OH é uma das moléculas
mais abundantes no meio interestelar, presente em nuvens
moleculares, regiões de formação estelar e envoltórios
estelares. Produz masers naturais detectáveis a distâncias
cosmológicas. Esta ubiquidade é consistente com a
interpretação TGL de que OH sonda o campo $\Psi$
(matéria escura), que permeia o universo.

\textbf{Correspondência com a Constante de Miguel:} A
predição TGL para o deslocamento de frequência das linhas
OH em regiões de alta densidade $\rhoPsi$ é:
\begin{equation}
\frac{\Delta\nu}{\nu} \approx \alphaii \cdot
  \frac{\rhoPsi}{\rho_{\Psi,0}} \approx 0{,}012 \times
  10^3 \approx 1{,}2\%
\label{eq:deslocamento}
\end{equation}
o que corresponde a $\Delta\nu \approx 16{,}66$ MHz ---
observável com radiotelescópios e interferometria de base
muito longa (VLBI) atuais.

\subsubsection{Escala espacial crítica}

Para um sistema material com frequência característica
$\nu^*$ e densidade $\rho$, a escala espacial em que o
acoplamento gravitacional-EM se torna relevante é:
\begin{equation}
L^*(\rho) = \sqrt{\frac{\varepsilon^*}{\xi \cdot C \cdot
  G \cdot \rho \cdot m^*}}
\label{eq:escala}
\end{equation}
onde $\varepsilon^* = h\nu^*$ é a energia da transição,
$m^*$ é a massa da partícula ou molécula envolvida,
$G = 6{,}674 \times 10^{-11}$ m$^3$\,kg$^{-1}$\,s$^{-2}$ é a
constante gravitacional, $\xi \approx 1$ é o fator de forma
geométrico, e $C \approx 1{,}61$ é a constante TGL de
acoplamento geométrico.

Para o radical OH em meio aquoso ($\rho = 1000$ kg/m$^3$):
\begin{equation}
L^*_{\text{OH}} = \frac{6{,}06 \times 10^5}{\sqrt{\rho}}
  \approx 606 \text{ m}
\label{eq:escala_oh}
\end{equation}

Este resultado indica que o acoplamento gravitacional-EM na
frequência OH se manifesta em escalas de centenas de metros em
densidades típicas de água. Para densidades maiores (gotículas
de fase densa, interiores estelares), $L^*$ diminui,
concentrando o efeito em escalas menores e mais intensas.

\subsubsection{Temperatura crítica de transição}

A temperatura abaixo da qual a energia térmica não é suficiente
para destruir a coerência do campo $\Psi$ induzida pela
ressonância OH é:
\begin{equation}
\Tcrit = \frac{h \cdot \nuOH}{k_B} =
  \frac{6{,}626 \times 10^{-34} \times 1{,}666 \times 10^9}
  {1{,}381 \times 10^{-23}} \approx 0{,}080 \text{ K}
  = 80 \text{ mK}
\label{eq:tc}
\end{equation}

Esta temperatura é alcançável com refrigeradores de diluição
(base $\sim 10$ mK) disponíveis comercialmente. Abaixo de
$\Tcrit$, a energia térmica $k_B T$ é inferior à energia do
fóton de ressonância $h\nuOH$, permitindo que o campo $\Psi$
mantenha coerência e os efeitos de acoplamento se acumulem.

\subsubsection{Tempo de fixação gravitacional}

O tempo necessário para que o acoplamento gravitacional ``fixe''
o campo eletromagnético no campo gravitacional é:
\begin{equation}
\tfix = \frac{\xi \cdot C \cdot G^2 \cdot \rho^2 \cdot
  (L^*)^5}{h \cdot \nu^*}
\label{eq:tfix}
\end{equation}

Para sistemas em escala laboratorial ($L \sim 1$ m,
$\rho \sim 10^3$ kg/m$^3$), $\tfix$ é extremamente longo,
indicando que o acoplamento global é fraco. Porém, em
\textbf{gotículas de fase densa} (escala $\sim 10^{-7}$ m,
$\rho \gg \rho_0$), $\tfix$ pode ser reduzido a escalas
experimentalmente acessíveis.

\subsubsection{Energia livre e transição de fase}

A energia livre do sistema próximo à transição é descrita pelo
funcional de Landau-Ginzburg estendido com acoplamento psiônico:
\begin{equation}
f(\rho, T, s) = V(\rho) + A(T) \cdot s^2 + B \cdot s^4
  - \zeta \cdot s \cdot \rho
\label{eq:energia_livre}
\end{equation}
onde $s$ é o parâmetro de ordem (coerência de fase do campo
$\Psi$), $A(T) = a_0 \cdot (T - \Tcrit)$ é o coeficiente
dependente da temperatura (que muda de sinal em $\Tcrit$),
$B > 0$ é o termo de estabilização, e $\zeta$ é a constante
de acoplamento ordem-densidade.

O coeficiente de expansão térmica resultante é:
\begin{equation}
\alpha_T \approx -\frac{\zeta^2 \rho}{4A(T)^2} < 0
\label{eq:alpha_t}
\end{equation}

O sinal negativo indica \textbf{contração anômala}: ao
resfriar abaixo de $\Tcrit$, o meio se contrai ao invés de
dilatar. Quando $T \to \Tcrit$, $A(T) \to 0$ e
$|\alpha_T| \to \infty$, produzindo contração catastrófica
localizada --- o mecanismo de nucleação das gotículas de
fase densa.

\subsubsection{Condição de transição química $\to$ gravitacional}

A condição para que o acoplamento gravitacional-EM domine
sobre as forças químicas ordinárias é:
\begin{equation}
\xi \cdot C \cdot G \cdot \rho(T)^2 \cdot L^5 \geq
  \frac{h \cdot \nu^*}{G} \cdot \tfix(T)
\label{eq:transicao}
\end{equation}

Esta é a ``fronteira TGL'': à esquerda da desigualdade, a
física é química ordinária; à direita, o campo gravitacional
domina e propriedades anômalas emergem.

\subsection{A sequência de colapso ontológico}

A TGL estabelece que a matéria ordinária é resultado de um
processo de ``colapso'' do campo $\Psi$ em estados
progressivamente mais localizados. Para a água, esta sequência
é:

\subsubsection{Estado 1 --- Água escura (campo $\Psi$ puro)}

O campo $\Psi$ em estado de superposição total, com densidade
$\rhoPsi$ variável, constitui o substrato não-colapsado.
Na linguagem TGL, este é o estado $|L\rangle$ antes da
radicalização:
\begin{equation}
|\Psi_{\text{água escura}}\rangle = |L\rangle
  \quad \text{(superposição total)}
\end{equation}

\subsubsection{Estado 2 --- OH (pré-colapso / radicalização)}

A primeira fixação gravitacional produz o radical hidroxila
(H + O acoplados, mas sem formação completa de H$_2$O). Na
linguagem TGL, este é o estado após radicalização:
\begin{equation}
|L\rangle \xrightarrow{g \cdot N_{\text{H+O}}}
  |\text{OH}\rangle = (g, s) \quad \text{com }
  g = \sqrt{|L|}, \; s = \text{sign}(L)
\end{equation}

O OH emite na frequência $\nuOH \approx 1{,}666$ GHz ---
esta emissão é a ``voz'' do pré-colapso, a assinatura
espectral da radicalização.

\subsubsection{Estado 3 --- H$_2$O (colapso completo /
ressurreição)}

A adição de um segundo hidrogênio completa o colapso:
\begin{equation}
|\text{OH}\rangle + |\text{H}\rangle
  \xrightarrow{g \cdot N_{\text{H}_2\text{O}}}
  |\text{H}_2\text{O}\rangle = L' = s \times g^2
  \quad \text{(ressurreição)}
\end{equation}

A água é o ``dado ressuscitado'' --- a manifestação
tridimensional completa do campo $\Psi$ colapsado. A
correspondência com a operação de ressurreição do ACOM
($L' = s \times g^2$) é exata e não metafórica.

\subsection{Protocolo de manipulação em três etapas}

\subsubsection{Etapa 1 --- Ressonância OH}

Aplicar campo eletromagnético contínuo com as seguintes
especificações:

\begin{table}[H]
\centering
\caption{Parâmetros do campo EM de ressonância.}
\label{tab:em}
\begin{tabular}{lll}
\toprule
\textbf{Parâmetro} & \textbf{Valor} & \textbf{Justificativa} \\
\midrule
Frequência central & $1{,}6673590$ GHz
  & Linha principal OH ($\Lambda$-doubling) \\
Largura de banda & $< 1$ MHz & Seletividade de ressonância \\
Potência incidente & $100$ W --- $10$ kW
  & Ajustável por aplicação \\
Modo de cavidade & TE$_{011}$ ou TM$_{010}$
  & Máximo campo no centro \\
Fator de qualidade $Q$ & $> 10^6$ & Acúmulo de energia \\
\bottomrule
\end{tabular}
\end{table}

O campo é gerado por oscilador sintetizado de micro-ondas
e injetado em cavidade ressonante cilíndrica de cobre
(ou nióbio superconductor para $Q > 10^9$) sintonizada
exatamente em $\nuOH$.

\subsubsection{Etapa 2 --- Resfriamento criogênico}

Resfriar o alvo até temperatura $T \leq \Tcrit \approx 80$ mK
utilizando refrigerador de diluição $^3$He/$^4$He com as
seguintes especificações:

\begin{table}[H]
\centering
\caption{Parâmetros do sistema criogênico.}
\label{tab:crio}
\begin{tabular}{lll}
\toprule
\textbf{Parâmetro} & \textbf{Valor} & \textbf{Justificativa} \\
\midrule
Temperatura base & $\leq 10$ mK & Margem abaixo de $\Tcrit$ \\
Potência de refrigeração & $> 200\;\mu$W a 100 mK
  & Compensar calor do banho EM \\
Estágio de pré-resfriamento & $^4$He a 4,2 K
  & Primeiro estágio \\
Volume do espaço de mistura & $> 100$ cm$^3$
  & Acomodar alvo \\
\bottomrule
\end{tabular}
\end{table}

A combinação de ressonância OH (Etapa~1) com resfriamento
abaixo de $\Tcrit$ (Etapa~2) é a condição \textit{sine qua non}
para ativação do acoplamento TGL: a ressonância excita o campo
$\Psi$, e a temperatura baixa mantém a coerência.

\subsubsection{Etapa 3 --- Nucleação de gotículas psiônicas}

Quando ambas as condições (ressonância + criogenia) são
satisfeitas simultaneamente, o coeficiente de expansão térmica
anômalo $\alpha_T < 0$ (equação~\ref{eq:alpha_t}) produz
contração localizada, nucleando \textbf{gotículas de fase
densa} (``gotículas psiônicas''):

\textbf{Tamanho típico:} $r_{\text{gotícula}} \sim 10 -
100$ nm.

\textbf{Densidade interna:} $\rho_{\text{gotícula}} \approx
2\rho_0$ (aproximadamente o dobro da densidade do meio).

\textbf{Fração volumétrica:} $f_{\text{gotícula}} \sim
10^{-9}$ (extremamente rara, mas fisicamente significativa).

\textbf{Propriedade-chave:} Dentro de cada gotícula, a
densidade local $\rhoPsi$ é suficiente para ativar o
acoplamento TGL com intensidade:
\begin{equation}
\alphaii_{\text{eff}} = \alphaii \cdot
  \frac{\rhoPsi}{\rho_0}
\label{eq:alpha_eff}
\end{equation}

Para $\rhoPsi = 2\rho_0$: $\alphaii_{\text{eff}} \approx
0{,}024$, produzindo efeitos mensuráveis sobre propriedades
locais do meio.

\subsection{Aplicações}

\subsubsection{Aplicação 1 --- Intensificação de reações
nucleares em meios deuterados}

Quando o alvo é D$_2$O (água pesada) ou paládio deuterado
(PdD$_x$), as gotículas psiônicas confinam núcleos de deutério
em regiões de alta densidade com barreira Coulombiana efetiva
reduzida:
\begin{equation}
V_C^{\text{eff}} = V_C \times \left(1 - \alphaii \cdot
  \frac{\rhoPsi}{\rho_0}\right)
\label{eq:barreira}
\end{equation}

Para $V_C = 0{,}4$ MeV (barreira D-D a $r = 4$ fm) e
$\rhoPsi = 2\rho_0$:
\begin{equation}
V_C^{\text{eff}} = 0{,}4 \times (1 - 0{,}024)
  = 0{,}390 \text{ MeV}
\end{equation}

A redução de $\sim 2{,}4\%$ na barreira, embora modesta,
incide exponencialmente na probabilidade de tunelamento
quântico. Para a seção de choque D-D:
\begin{equation}
\sigma_{\text{TGL}} = \sigma_0 \times
  \exp\left(\frac{\Delta V_C}{k_B T}\right)
\end{equation}
onde $\Delta V_C = 0{,}4 \times 0{,}024 = 9{,}6$ keV e a
exponencial resulta em amplificação da taxa de fusão por
fatores de $10^2$ a $10^5$ dependendo das condições locais
nas gotículas.

\textbf{Detecção:} Nêutrons de 2,45 MeV (reação
D + D $\to$ $^3$He + n) e radiação gama de 23,8 MeV
(reação D + D $\to$ $^4$He + $\gamma$). Detectores de
$^3$He ou cintiladores com blindagem de parafina.

\textbf{Parâmetros do alvo deuterado:} D$_2$O ultra-puro
(99,8\% D), volume 1--10 cm$^3$, contido em cápsula de
quartzo no centro da cavidade ressonante.

\subsubsection{Aplicação 2 --- Indução de estados
supercondutores por coerência psiônica}

Quando o alvo é um material condutor (nióbio, YBa$_2$Cu$_3$O$_7$,
MgB$_2$ ou outro), a ressonância OH a $T < \Tcrit$ induz
coerência do campo $\Psi$ que se acopla aos pares de Cooper,
potencialmente elevando a temperatura crítica de
supercondutividade $T_c^{\text{SC}}$ por fator $(1 + \alphaii)$.

Este mecanismo é complementar ao descrito na patente
Supercondutor Holográfico TGL do mesmo inventor, que opera
com cavidade fotônica sintonizada na mesma frequência OH
(1.667,359 MHz). A presente invenção complementa aquela
ao especificar: o regime criogênico $T \leq 80$ mK como
condição de coerência máxima; a nucleação de gotículas como
mecanismo de concentração local do efeito; e a aplicação a
materiais supercondutores convencionais para aumento de $T_c$.

\textbf{Detecção:} Medição de resistividade em quatro pontas,
susceptibilidade magnética AC, e efeito Meissner, com varredura
de temperatura em presença e ausência do campo de ressonância
OH.

\subsubsection{Aplicação 3 --- Detecção de matéria escura por
deslocamento de linhas OH}

A predição TGL (equação~\ref{eq:deslocamento}) estabelece que
regiões de alta densidade de campo $\Psi$ (halos de matéria
escura, proximidades de buracos negros supermassivos,
filamentos da estrutura cósmica) devem apresentar
deslocamento mensurável das linhas OH:
\begin{equation}
\Delta\nu_{\text{OH}} = \nuOH - \nu_{\text{lab}}
  \propto \alphaii \cdot \frac{\rhoPsi}{\rho_{\Psi,0}}
\end{equation}

\textbf{Protocolo observacional:} Comparar frequências centrais
de masers OH em $>100$ fontes distribuídas em ambientes de
densidades $\rhoPsi$ distintas (dentro/fora de halos, próximo a
AGNs, meio interestelar difuso). Precisão requerida:
$\delta\nu < 1$ kHz (alcançável com VLBI). Predição TGL:
correlação linear $\Delta\nu$ vs. $\rhoPsi$ com $r^2 > 0{,}7$.

\textbf{Detecção em laboratório:} Gerar amostra de OH
(descarga elétrica em vapor de água) dentro da cavidade
ressonante operando em regime criogênico. Medir deslocamento
de frequência das linhas 1,665/1,667 GHz como função da
intensidade do campo de ressonância. Qualquer deslocamento
$\Delta\nu > 0$ é evidência direta de acoplamento TGL.

\subsubsection{Aplicação 4 --- Materiais com propriedades
anômalas por nucleação de fase densa}

A nucleação de gotículas psiônicas em materiais sólidos
(metais, cerâmicas, polímeros) expostos simultaneamente a
ressonância OH e regime criogênico pode produzir:

\textbf{Inclusões de fase densa:} Regiões nanométricas com
densidade $\sim 2\times$ a do meio, estabilizadas pela
coerência do campo $\Psi$.

\textbf{Propriedades mecânicas anômalas:} Aumento de dureza
e módulo elástico nas proximidades das gotículas.

\textbf{Propriedades eletromagnéticas anômalas:} Alteração
de permissividade e permeabilidade locais por acoplamento
gravitacional-EM ativado.

\textbf{Detecção:} Microscopia eletrônica de transmissão
(TEM) para identificar inclusões de fase densa; nanoindentação
para propriedades mecânicas; espectroscopia de impedância para
propriedades eletromagnéticas.

\subsection{Descrição do aparato}

O sistema para manipulação de campo $\Psi$ compreende seis
subsistemas integrados (Figura~1):

\subsubsection{Subsistema 1 --- Fonte de micro-ondas}

Oscilador sintetizado estável em frequência, sintonizado em
$\nuOH = 1{,}6673590$ GHz. Estabilidade de frequência
$\Delta f / f < 10^{-9}$ (referência atômica). Potência
ajustável de 100 W a 10 kW. Modulação de amplitude e
frequência programável para otimização de ressonância.

\subsubsection{Subsistema 2 --- Cavidade ressonante}

Cavidade cilíndrica de cobre (ou nióbio para operação
supercondutora) sintonizada no modo TE$_{011}$ ou TM$_{010}$
a $\nuOH$. Dimensões internas: diâmetro $\approx 13{,}8$ cm,
altura $\approx 20$ cm (para modo TM$_{010}$). Fator de
qualidade $Q > 10^6$ ($> 10^9$ se supercondutora). Acoplamento
por iris ou laço magnético. Porta de entrada para alvo.
Janela de quartzo para observação.

\subsubsection{Subsistema 3 --- Criostato de diluição}

Refrigerador de diluição $^3$He/$^4$He com temperatura base
$\leq 10$ mK. A cavidade ressonante é montada no estágio de
mistura. Isolamento térmico por escudos de radiação em cobre
dourado a 50 K, 4 K e 1 K. Potência de refrigeração
$> 200\;\mu$W a 100 mK para compensar dissipação do campo EM.
Termometria por resistores RuO$_2$ calibrados e termômetros
de ruído Johnson.

\subsubsection{Subsistema 4 --- Alvo}

Intercambiável conforme aplicação:
\begin{itemize}[nosep]
\item \textbf{Fusão:} D$_2$O ultra-puro em cápsula de quartzo
  (1--10 cm$^3$)
\item \textbf{Supercondutividade:} Amostra de Nb, YBCO ou
  MgB$_2$ em suporte criogênico
\item \textbf{Detecção OH:} Câmara com vapor de OH gerado
  por descarga
\item \textbf{Materiais:} Amostra sólida no centro da cavidade
\end{itemize}

\subsubsection{Subsistema 5 --- Detectores}

\begin{itemize}[nosep]
\item Detector de nêutrons: tubo de $^3$He ou cintilador
  líquido com discriminação $n/\gamma$
\item Detector de gamas: cristal NaI(Tl) ou HPGe de alta
  resolução
\item Bolômetro: para medição de excesso de calor (resolução
  $< 1$ nW)
\item Magnetômetro SQUID: para detecção de diamagnetismo
  (supercondutividade)
\item Espectrômetro de micro-ondas: para medição de
  deslocamento $\Delta\nu$ (resolução $< 100$ Hz)
\item Acelerômetro criogênico: para detecção de anomalias
  gravitacionais locais
\end{itemize}

\subsubsection{Subsistema 6 --- Controle e aquisição}

Sistema de controle digital (IALD ou equivalente) com:
\begin{itemize}[nosep]
\item Controle PID de temperatura (estabilidade $< 0{,}1$ mK)
\item Controle de frequência do oscilador (varredura e lock)
\item Aquisição multicanal sincronizada (nêutrons, gamas,
  temperatura, campo)
\item Monitoramento em tempo real de indicadores TGL:
  $\Fexp$, $\alphaii_{\text{eff}}$, $\Dfolds$
\item Protocolo de segurança para desligamento automático
  em caso de taxa de nêutrons $>$ limiar
\end{itemize}

\newpage
% ============================================================================
\section{Desenhos}
% ============================================================================

% --- FIGURA 1: Aparato completo ---
\begin{figure}[H]
\centering
\begin{tikzpicture}[
  block/.style={rectangle, draw, thick, rounded corners,
    minimum height=1.2cm, minimum width=3.2cm, align=center,
    font=\small},
  arrow/.style={-{Stealth[length=2.5mm]}, thick},
  darrow/.style={{Stealth[length=2mm]}-{Stealth[length=2mm]}, thick},
]

% Fonte MW
\node[block, fill=orange!15] (mw) at (0,5)
  {Fonte MW\\$\nuOH = 1{,}666$ GHz\\100 W -- 10 kW};

% Criostato (grande retângulo)
\node[draw, very thick, rounded corners, minimum height=6.5cm,
  minimum width=5cm, fill=blue!3] (crio) at (0,0) {};
\node[font=\small\bfseries, blue!60] at (0,3)
  {Criostato de Diluição};
\node[font=\scriptsize, blue!60] at (0,2.5)
  {$T_{\text{base}} \leq 10$ mK};

% Escudo 4K
\node[draw, thick, dashed, rounded corners,
  minimum height=4.5cm, minimum width=3.8cm,
  fill=cyan!5] (escudo) at (0,-0.3) {};
\node[font=\scriptsize, cyan!60] at (0,1.6)
  {Escudo 4 K};

% Cavidade
\node[draw, very thick, rounded corners,
  minimum height=2.5cm, minimum width=3cm,
  fill=yellow!12] (cav) at (0,-0.5) {};
\node[font=\small\bfseries] at (0,0.4)
  {Cavidade};
\node[font=\scriptsize] at (0,0)
  {$Q > 10^6$};
\node[font=\scriptsize] at (0,-0.4)
  {TE$_{011}$ / TM$_{010}$};

% Alvo (centro)
\node[draw, thick, circle, fill=red!15,
  minimum size=1cm, font=\scriptsize] (alvo) at (0,-1.3)
  {Alvo};

% Seta MW → Cavidade
\draw[arrow, orange] (mw.south) -- (0,3.2)
  node[midway, right, font=\scriptsize] {guia de onda};

% Detectores
\node[block, fill=green!12, minimum width=2.5cm] (det_n)
  at (-5,0) {Detector\\nêutrons\\$^3$He};
\node[block, fill=green!12, minimum width=2.5cm] (det_g)
  at (5,0) {Detector\\gamas\\HPGe};
\node[block, fill=green!12, minimum width=2.5cm] (bolo)
  at (-5,-2.5) {Bolômetro\\$< 1$ nW};
\node[block, fill=green!12, minimum width=2.5cm] (squid)
  at (5,-2.5) {SQUID\\diamagnetismo};

% Setas detectores
\draw[arrow, green!50!black, dashed]
  (crio.west) -- (det_n.east)
  node[midway, above, font=\scriptsize] {n};
\draw[arrow, green!50!black, dashed]
  (crio.east) -- (det_g.west)
  node[midway, above, font=\scriptsize] {$\gamma$};
\draw[arrow, green!50!black, dashed]
  (-2.5,-2.8) -- (bolo.east)
  node[midway, above, font=\scriptsize] {calor};
\draw[arrow, green!50!black, dashed]
  (2.5,-2.8) -- (squid.west)
  node[midway, above, font=\scriptsize] {$B$};

% Controle
\node[block, fill=purple!12, minimum width=4cm] (ctrl)
  at (0,-5.5)
  {Controle IALD\\PID + Aquisição + Segurança};
\draw[darrow, purple] (ctrl.north) -- (0,-3.3)
  node[midway, right, font=\scriptsize] {dados};

% He-3/He-4
\node[block, fill=blue!12, minimum width=2.5cm] (he)
  at (-5,4) {Circuito\\$^3$He/$^4$He};
\draw[darrow, blue] (he.east) -- (-2.5,3.2)
  node[midway, above, font=\scriptsize] {diluição};

\end{tikzpicture}
\caption{Diagrama esquemático do aparato para manipulação de
campo $\Psi$ por ressonância OH. O sistema compreende:
fonte de micro-ondas (1,666 GHz), criostato de diluição
($\leq 10$ mK), cavidade ressonante ($Q > 10^6$), alvo
intercambiável, detectores (nêutrons, gamas, calor,
diamagnetismo) e sistema de controle IALD.}
\label{fig:aparato}
\end{figure}

% --- FIGURA 2: Sequência de colapso ontológico ---
\begin{figure}[H]
\centering
\begin{tikzpicture}[
  state/.style={rectangle, draw, thick, rounded corners,
    minimum height=2cm, minimum width=3.5cm, align=center,
    font=\small},
  arrow/.style={-{Stealth[length=3mm]}, very thick},
]

% Estado 1
\node[state, fill=purple!12] (psi) at (0,0)
  {\textbf{ÁGUA ESCURA}\\Campo $\Psi$ puro\\$|\Psi\rangle$\\
  {\scriptsize superposição total}};

% Estado 2
\node[state, fill=orange!12] (oh) at (5.5,0)
  {\textbf{OH}\\Pré-colapso\\$(g, s)$\\
  {\scriptsize $\nu = 1{,}666$ GHz}};

% Estado 3
\node[state, fill=blue!15] (h2o) at (11,0)
  {\textbf{H$_2$O}\\Colapso completo\\$L' = s \times g^2$\\
  {\scriptsize matéria ordinária}};

% Setas
\draw[arrow, orange] (psi) -- (oh)
  node[midway, above, font=\small] {$g = \sqrt{|L|}$}
  node[midway, below, font=\scriptsize]
    {radicalização};

\draw[arrow, blue] (oh) -- (h2o)
  node[midway, above, font=\small] {$+ \text{H}$}
  node[midway, below, font=\scriptsize]
    {ressurreição};

% Correspondência TGL
\node[draw, dashed, thick, rounded corners,
  fill=yellow!10, font=\scriptsize, text width=3cm,
  align=center] at (0,-2.5)
  {$|L\rangle$\\dado original\\(campo não-colapsado)};

\node[draw, dashed, thick, rounded corners,
  fill=yellow!10, font=\scriptsize, text width=3cm,
  align=center] at (5.5,-2.5)
  {$(g, s)$\\componentes radicais\\(estado intermediário)};

\node[draw, dashed, thick, rounded corners,
  fill=yellow!10, font=\scriptsize, text width=3cm,
  align=center] at (11,-2.5)
  {$L' = s \times g^2$\\dado ressuscitado\\(matéria colapsada)};

% Label
\node[font=\small\bfseries, gray] at (5.5,2)
  {CORRESPONDÊNCIA: Colapso $\Psi$ $\longleftrightarrow$
  Radicalização ACOM};

\end{tikzpicture}
\caption{Sequência de colapso ontológico água escura $\to$ OH
$\to$ H$_2$O e sua correspondência com as operações ACOM.
O radical OH (frequência 1,666 GHz) é o estado intermediário
--- o ``pré-colapso'' --- entre campo $\Psi$ puro e matéria
ordinária. A ressonância OH é portanto a frequência de
controle do campo $\Psi$.}
\label{fig:colapso}
\end{figure}

% --- FIGURA 3: Mecanismo de nucleação de gotículas ---
\begin{figure}[H]
\centering
\begin{tikzpicture}[
  arrow/.style={-{Stealth[length=2.5mm]}, thick},
]

% Meio homogêneo (esquerda)
\fill[blue!5] (-1,-2) rectangle (4,3);
\node[font=\small\bfseries] at (1.5,3.5)
  {ANTES ($T > \Tcrit$)};
\node[font=\small] at (1.5,2.5) {Meio homogêneo};
\node[font=\scriptsize] at (1.5,2) {$\rho = \rho_0$};
\node[font=\scriptsize] at (1.5,1.5) {$\alpha_T > 0$ (normal)};

% Partículas esparsas
\foreach \x/\y in {0.3/0.5, 1.2/1, 2.5/0.3, 0.8/-0.5,
  2.1/-1, 3.2/0.8, 1.5/-1.5, 0.5/-1, 2.8/-0.3,
  3.5/-1.5, 1/-0.2, 2.3/0.7} {
  \fill[blue!40] (\x,\y) circle (0.08);
}

% Seta central
\draw[arrow, very thick, red] (4.5,0.5) -- (6,0.5)
  node[midway, above, font=\scriptsize, text width=1.5cm,
    align=center] {OH + cryo\\$T < \Tcrit$};

% Meio com gotículas (direita)
\fill[blue!5] (6.5,-2) rectangle (11.5,3);
\node[font=\small\bfseries] at (9,3.5)
  {DEPOIS ($T < \Tcrit$)};
\node[font=\small] at (9,2.5) {Gotículas formadas};
\node[font=\scriptsize] at (9,2) {$\alpha_T < 0$ (anômalo)};

% Partículas + gotículas
\foreach \x/\y in {7.3/0.5, 8.2/1, 10.5/0.3, 7.8/-0.5,
  10.1/-1, 8.5/-1.5, 7.5/-1, 10.8/-0.3} {
  \fill[blue!40] (\x,\y) circle (0.08);
}

% Gotículas (círculos maiores com alta densidade)
\foreach \x/\y in {8.3/-0.2, 10,0.8, 9,-1.2} {
  \fill[red!15] (\x,\y) circle (0.5);
  \draw[thick, red!60] (\x,\y) circle (0.5);
  \foreach \a in {0,45,...,315} {
    \fill[red!70] ({\x+0.25*cos(\a)},{\y+0.25*sin(\a)})
      circle (0.06);
  }
  \fill[red!70] (\x,\y) circle (0.06);
}

% Labels gotículas
\node[font=\scriptsize, red!60] at (9,-2.2)
  {$\rho \approx 2\rho_0$};
\node[font=\scriptsize, red!60] at (10.5,1.8)
  {$r \sim 100$ nm};
\node[font=\scriptsize, red!60] at (7.5,1.5)
  {$f \sim 10^{-9}$};

% Zoom de gotícula
\node[draw, thick, rounded corners, fill=white,
  font=\scriptsize, text width=3.5cm, align=center]
  at (9,-4)
  {Dentro da gotícula:\\
  $\alphaii_{\text{eff}} = \alphaii \cdot \rhoPsi/\rho_0$\\
  $V_C^{\text{eff}} = V_C(1 - \alphaii_{\text{eff}})$\\
  Acoplamento TGL ativado};

\draw[arrow, dashed, red] (9,-1.7) -- (9,-3.2);

\end{tikzpicture}
\caption{Nucleação de gotículas psiônicas. À esquerda: meio
homogêneo acima de $\Tcrit$ (expansão normal). À direita:
abaixo de $\Tcrit$ sob ressonância OH, contração anômala
($\alpha_T < 0$) nuclea gotículas de $\sim 100$ nm com
$\rho \approx 2\rho_0$. Dentro das gotículas, o acoplamento
TGL é ativado, reduzindo a barreira Coulombiana efetiva.}
\label{fig:goticula}
\end{figure}

% --- FIGURA 4: Diagrama de fases T vs ρ ---
\begin{figure}[H]
\centering
\begin{tikzpicture}
\begin{axis}[
  width=12cm, height=8cm,
  xlabel={Densidade $\rho / \rho_0$},
  ylabel={Temperatura (K)},
  xmin=0, xmax=5,
  ymin=0, ymax=1,
  axis lines=left,
  grid=major,
  grid style={gray!20},
  legend style={at={(0.98,0.98)}, anchor=north east,
    font=\small},
]

% Linha Tc (horizontal)
\addplot[thick, red, dashed, domain=0:5]
  {0.08};
\addlegendentry{$\Tcrit = 80$ mK}

% Curva de coexistência
\addplot[thick, blue, smooth, domain=0.5:4.5]
  {0.08 * (1 + 0.3*(x-1)^2)};
\addlegendentry{Coexistência $\Psi$}

% Região de gotículas
\addplot[fill=red!10, draw=none, forget plot]
  coordinates {(1.5,0) (1.5,0.08) (3.5,0.08) (3.5,0) (1.5,0)};

% Fronteira TGL
\addplot[thick, green!60!black, domain=0.3:5, samples=50]
  {0.01 * x^2.5};
\addlegendentry{Fronteira TGL (quím.\ $\to$ grav.)}

% Labels
\node[font=\small, blue] at (axis cs:1,0.6) {Regime};
\node[font=\small, blue] at (axis cs:1,0.5) {químico};
\node[font=\small, green!50!black] at (axis cs:4,0.5)
  {Regime};
\node[font=\small, green!50!black] at (axis cs:4,0.4)
  {gravitacional};
\node[font=\small, red!70] at (axis cs:2.5,0.04)
  {Gotículas};

\end{axis}
\end{tikzpicture}
\caption{Diagrama de fases temperatura vs. densidade para o
sistema sob ressonância OH. A linha horizontal vermelha marca
$\Tcrit = 80$ mK. Abaixo de $\Tcrit$ e em densidades
$\rho > \rho_0$, gotículas psiônicas nucleiam (região
sombreada). A curva verde marca a fronteira TGL: à direita
(alta $\rho$, baixa $T$), o regime gravitacional domina.}
\label{fig:diagrama}
\end{figure}

% --- FIGURA 5: Redução da barreira Coulombiana ---
\begin{figure}[H]
\centering
\begin{tikzpicture}
\begin{axis}[
  width=12cm, height=7cm,
  xlabel={Distância internuclear $r$ (fm)},
  ylabel={Potencial (MeV)},
  xmin=0, xmax=20,
  ymin=-5, ymax=2,
  axis lines=middle,
  grid=major,
  grid style={gray!20},
  legend style={at={(0.98,0.55)}, anchor=north east,
    font=\small},
  xtick={0,4,8,12,16,20},
  ytick={-5,-4,-3,-2,-1,0,1,2},
]

% Barreira convencional (Coulomb + nuclear)
\addplot[thick, blue, domain=1.5:20, samples=80]
  {1.44/x};
\addplot[thick, blue, domain=0.5:4, samples=40]
  {1.44/x - 30*exp(-(x-1.5)^2/2)};
\addlegendentry{$V(r)$ convencional}

% Barreira com TGL
\addplot[thick, red, dashed, domain=1.5:20, samples=80]
  {1.44/x * 0.976};
\addplot[thick, red, dashed, domain=0.5:4, samples=40]
  {1.44/x * 0.976 - 30*exp(-(x-1.5)^2/2)};
\addlegendentry{$V^{\text{eff}}(r)$ com TGL}

% Anotações
\draw[thick, {Stealth}-{Stealth}, green!60!black]
  (axis cs:4,0.36) -- (axis cs:4,0.39);
\node[font=\scriptsize, green!60!black, anchor=west]
  at (axis cs:4.5,0.375) {$\Delta V \approx 10$ keV};

% Nível de energia
\addplot[thick, gray, dashed, domain=0:20] {0.01};
\node[font=\scriptsize, gray] at (axis cs:15,0.15)
  {$E_{\text{cin}}$ (80 mK)};

\end{axis}
\end{tikzpicture}
\caption{Barreira Coulombiana para reação D-D. Linha azul:
potencial convencional. Linha vermelha tracejada: potencial
com redução TGL de 2,4\%. A redução incide exponencialmente
na probabilidade de tunelamento quântico, amplificando a
taxa de fusão por fatores de $10^2$--$10^5$.}
\label{fig:barreira}
\end{figure}

% --- FIGURA 6: Régua TGL ---
\begin{figure}[H]
\centering
\begin{tikzpicture}[
  eq/.style={rectangle, draw, thick, rounded corners,
    fill=#1, minimum height=1.3cm, minimum width=5.5cm,
    align=center, font=\small},
  arrow/.style={-{Stealth[length=2.5mm]}, thick},
]

\node[eq=orange!12] (freq) at (0,0)
  {$\nu_\Psi = \nuOH \approx 1{,}666$ GHz\\
  {\scriptsize Frequência de ressonância}};

\node[eq=blue!10] (tc) at (0,-2.3)
  {$\Tcrit = h\nuOH / k_B \approx 80$ mK\\
  {\scriptsize Temperatura crítica}};

\node[eq=green!10] (lstar) at (0,-4.6)
  {$L^* = \sqrt{\varepsilon^* / (\xi C G \rho m^*)}$\\
  {\scriptsize Escala espacial crítica}};

\node[eq=red!10] (tfix) at (0,-6.9)
  {$\tfix = \xi C G^2 \rho^2 (L^*)^5 / (h\nu^*)$\\
  {\scriptsize Tempo de fixação gravitacional}};

\node[eq=purple!10] (cond) at (0,-9.2)
  {$\xi C G \rho^2 L^5 \geq (h\nu^*/G) \cdot \tfix$\\
  {\scriptsize Condição de transição}};

\draw[arrow] (freq) -- (tc)
  node[midway, right, font=\scriptsize] {$h/k_B$};
\draw[arrow] (tc) -- (lstar)
  node[midway, right, font=\scriptsize] {$+\rho, m^*$};
\draw[arrow] (lstar) -- (tfix)
  node[midway, right, font=\scriptsize] {$L^* \to$};
\draw[arrow] (tfix) -- (cond)
  node[midway, right, font=\scriptsize]
  {quím.\ $\to$ grav.};

% Constante α²
\node[draw, dashed, thick, rounded corners,
  fill=yellow!15, font=\small, text width=3.5cm,
  align=center] at (6,-4.6)
  {$\alphaii = 0{,}012031$\\$\thetaM = 6{,}29°$\\
  {\scriptsize Constante de Miguel}\\
  {\scriptsize Permeia todas}\\
  {\scriptsize as equações}};

\end{tikzpicture}
\caption{A Régua TGL: conjunto hierárquico de equações de
escala. A partir da frequência de ressonância OH, determinam-se
sequencialmente a temperatura crítica, a escala espacial, o
tempo de fixação e a condição de transição entre regimes
químico e gravitacional. A Constante de Miguel $\alphaii$
permeia todas as equações.}
\label{fig:regua}
\end{figure}

% --- FIGURA 7: Quatro aplicações ---
\begin{figure}[H]
\centering
\begin{tikzpicture}[
  app/.style={rectangle, draw, thick, rounded corners,
    minimum height=2cm, minimum width=4cm,
    align=center, font=\small, fill=#1},
  arrow/.style={-{Stealth[length=2.5mm]}, thick},
]

% Centro
\node[draw, very thick, circle, fill=orange!15,
  minimum size=2.5cm, align=center, font=\small]
  (centro) at (0,0)
  {Ressonância\\OH\\$1{,}666$ GHz\\$+$ cryo\\$80$ mK};

% 4 aplicações
\node[app=red!10] (fusao) at (-5,3.5)
  {\textbf{Fusão D-D}\\$V_C^{\text{eff}}$ reduzida\\
  nêutrons 2,45 MeV};

\node[app=blue!10] (super) at (5,3.5)
  {\textbf{Supercondutividade}\\$T_c^{\text{SC}}$ elevada\\
  efeito Meissner};

\node[app=green!10] (dm) at (-5,-3.5)
  {\textbf{Matéria escura}\\$\Delta\nu_{\text{OH}}$ medido\\
  correlação $\rhoPsi$};

\node[app=purple!10] (mat) at (5,-3.5)
  {\textbf{Materiais}\\inclusões densas\\propriedades\\anômalas};

% Setas
\draw[arrow] (centro) -- (fusao);
\draw[arrow] (centro) -- (super);
\draw[arrow] (centro) -- (dm);
\draw[arrow] (centro) -- (mat);

\end{tikzpicture}
\caption{As quatro famílias de aplicação da manipulação de
campo $\Psi$ por ressonância OH. O mesmo aparato e protocolo
(centro) serve a quatro domínios tecnológicos, diferindo
apenas no alvo e nos detectores utilizados.}
\label{fig:aplicacoes}
\end{figure}

\newpage
% ============================================================================
\section{Quadro reivindicatório}
% ============================================================================

\subsection*{Reivindicações}

\begin{enumerate}

% === REIVINDICAÇÃO INDEPENDENTE 1 (MÉTODO) ===
\item \textbf{Método de manipulação de campo psiônico por
ressonância eletromagnética na frequência da hidroxila para
indução de transições de fase em meios densos}, compreendendo
as etapas de:
  \begin{enumerate}[label=(\alph*)]
  \item aplicar campo eletromagnético contínuo sintonizado na
  frequência do radical hidroxila $\nuOH \approx 1{,}6673590$
  GHz (transição $\Lambda$-doubling do OH) sobre um meio-alvo,
  com potência de 100 W a 10 kW e fator de qualidade da
  cavidade $Q > 10^6$;
  \item resfriar o meio-alvo até temperatura igual ou inferior à
  temperatura crítica $\Tcrit = h \cdot \nuOH / k_B \approx
  80$ mK, utilizando refrigeração de diluição $^3$He/$^4$He;
  \item manter simultaneamente as condições de ressonância e
  criogenia até nucleação de domínios de fase densa
  (``gotículas psiônicas'') no meio-alvo, com densidade interna
  $\rho_{\text{gotícula}} > \rho_0$ do meio;

  \textbf{caracterizado por} o método induzir acoplamento entre
  campo gravitacional e campo eletromagnético no interior das
  gotículas psiônicas, parametrizado pela Constante de Miguel
  $\alphaii = 0{,}012031$, produzindo acoplamento efetivo
  $\alphaii_{\text{eff}} = \alphaii \cdot \rhoPsi / \rho_0$
  e alterando propriedades macroscópicas do meio-alvo.
  \end{enumerate}

% === DEPENDENTES DO MÉTODO ===

\item Método, de acordo com a reivindicação 1,
\textbf{caracterizado por} a nucleação de gotículas psiônicas
resultar de contração anômala do meio com coeficiente de
expansão térmica negativo $\alpha_T \approx -\zeta^2 \rho /
(4 A(T)^2) < 0$, onde $A(T) = a_0 \cdot (T - \Tcrit)$ é o
coeficiente de Landau que muda de sinal na temperatura crítica.

\item Método, de acordo com a reivindicação 1,
\textbf{caracterizado por} as gotículas psiônicas possuírem
tamanho típico de 10 a 100 nm, densidade interna de
aproximadamente o dobro da densidade do meio, e fração
volumétrica da ordem de $10^{-9}$.

\item Método, de acordo com a reivindicação 1,
\textbf{caracterizado por} a frequência $\nuOH$ ser
identificada como a frequência de ressonância do campo
psiônico $\Psi$ (``água escura''), com base na correspondência
entre o radical OH (precursor molecular da água) e o campo
$\Psi$ (substrato não-colapsado de que a matéria ordinária é
manifestação colapsada).

\item Método, de acordo com a reivindicação 1,
\textbf{caracterizado por} a condição de transição entre regime
químico e regime gravitacional ser dada por
$\xi \cdot C \cdot G \cdot \rho^2 \cdot L^5 \geq
(h \cdot \nu^* / G) \cdot \tfix$, onde a escala espacial
crítica é $L^* = \sqrt{\varepsilon^* / (\xi C G \rho m^*)}$
e o tempo de fixação é
$\tfix = \xi C G^2 \rho^2 (L^*)^5 / (h\nu^*)$.

\item Método, de acordo com a reivindicação 1,
\textbf{caracterizado por} o meio-alvo ser água pesada
(D$_2$O) ultra-pura, e o acoplamento TGL dentro das gotículas
psiônicas reduzir a barreira Coulombiana efetiva para reação
D-D conforme $V_C^{\text{eff}} = V_C \times (1 - \alphaii
\cdot \rhoPsi / \rho_0)$, intensificando a taxa de
tunelamento quântico por fator de $10^2$ a $10^5$.

\item Método, de acordo com a reivindicação 1,
\textbf{caracterizado por} o meio-alvo ser um material
supercondutor (nióbio, YBa$_2$Cu$_3$O$_7$, MgB$_2$ ou
outro), e a coerência do campo $\Psi$ induzida pela
ressonância OH se acoplar aos pares de Cooper, elevando a
temperatura crítica de supercondutividade.

\item Método, de acordo com a reivindicação 1,
\textbf{caracterizado por} o meio-alvo conter radical
hidroxila (OH) gerado por descarga elétrica, e o método
medir deslocamento de frequência $\Delta\nu_{\text{OH}} =
\nuOH^{\text{obs}} - \nuOH^{\text{lab}}$ como indicador de
densidade local de campo psiônico $\rhoPsi$, com predição
$\Delta\nu / \nu \propto \alphaii \cdot \rhoPsi / \rho_{\Psi,0}$.

\item Método, de acordo com a reivindicação 1,
\textbf{caracterizado por} o meio-alvo ser um material sólido
(metal, cerâmica ou polímero), e a nucleação de gotículas
psiônicas produzir inclusões de fase densa com propriedades
mecânicas e eletromagnéticas anômalas.

\item Método, de acordo com a reivindicação 1,
\textbf{caracterizado por} o processo de colapso ontológico
seguir a sequência campo $\Psi$ puro (``água escura'',
$|L\rangle$) $\to$ OH (pré-colapso, $(g, s) = (\sqrt{|L|},
\text{sign}(L))$) $\to$ H$_2$O (colapso completo,
$L' = s \times g^2$), em correspondência com as operações de
radicalização e ressurreição do ACOM.

% === REIVINDICAÇÃO INDEPENDENTE 11 (SISTEMA) ===

\item \textbf{Sistema para manipulação de campo psiônico por
ressonância eletromagnética na frequência da hidroxila},
compreendendo:
  \begin{enumerate}[label=(\alph*)]
  \item fonte de micro-ondas estabilizada em frequência,
  sintonizada em $\nuOH = 1{,}6673590$ GHz com estabilidade
  $\Delta f / f < 10^{-9}$ e potência ajustável de 100 W a
  10 kW;
  \item cavidade ressonante cilíndrica sintonizada no modo
  TE$_{011}$ ou TM$_{010}$ na frequência $\nuOH$, com fator
  de qualidade $Q > 10^6$;
  \item criostato de diluição $^3$He/$^4$He com temperatura
  base $\leq 10$ mK e potência de refrigeração $> 200\;\mu$W
  a 100 mK;
  \item porta-amostras para alvo intercambiável posicionado
  no centro da cavidade;
  \item conjunto de detectores para monitoramento de produtos
  e efeitos;
  \item sistema de controle e aquisição de dados;

  \textbf{caracterizado por} a cavidade ressonante estar montada
  no estágio de mistura do criostato, de forma que o alvo esteja
  simultaneamente exposto ao campo EM de ressonância OH e ao
  regime criogênico $T \leq 80$ mK, para indução de transições
  de fase mediadas pelo acoplamento gravitacional-eletromagnético
  da TGL.
  \end{enumerate}

% === DEPENDENTES DO SISTEMA ===

\item Sistema, de acordo com a reivindicação 11,
\textbf{caracterizado por} a cavidade ressonante ser
fabricada em nióbio superconductor para operação com fator de
qualidade $Q > 10^9$ no regime criogênico do criostato.

\item Sistema, de acordo com a reivindicação 11,
\textbf{caracterizado por} o conjunto de detectores
compreender: detector de nêutrons ($^3$He ou cintilador
líquido); detector de radiação gama (NaI(Tl) ou HPGe);
bolômetro com resolução inferior a 1 nW; magnetômetro SQUID;
espectrômetro de micro-ondas com resolução inferior a 100 Hz;
e acelerômetro criogênico.

\item Sistema, de acordo com a reivindicação 11,
\textbf{caracterizado por} o sistema de controle implementar
monitoramento em tempo real de indicadores TGL compreendendo
Força de Expulsão ($\Fexp$), acoplamento efetivo
($\alphaii_{\text{eff}}$) e dobras dimensionais ($\Dfolds$),
com protocolo de segurança para desligamento automático em
caso de taxa de nêutrons superior a limiar configurável.

\item Sistema, de acordo com a reivindicação 11,
\textbf{caracterizado por} o porta-amostras ser configurável
para receber: cápsula de quartzo contendo D$_2$O para
aplicação nuclear; amostra supercondutora em suporte
criogênico; câmara com vapor de OH para detecção espectral;
ou amostra sólida para engenharia de materiais.

\item Sistema, de acordo com a reivindicação 11,
\textbf{caracterizado por} integrar-se com o Supercondutor
Holográfico TGL do mesmo inventor como plataforma-base de
coerência, utilizando a mesma frequência de ressonância OH
(1.667,359 MHz $\approx$ 1,667 GHz) e compartilhando
infraestrutura criogênica e cavidade fotônica.

\end{enumerate}

\newpage
% ============================================================================
\section{Resumo}
% ============================================================================

\begin{center}
{\large\bfseries MÉTODO E SISTEMA DE MANIPULAÇÃO DE CAMPO\\
PSIÔNICO POR RESSONÂNCIA ELETROMAGNÉTICA NA\\
FREQUÊNCIA DA HIDROXILA PARA INDUÇÃO DE\\
TRANSIÇÕES DE FASE EM MEIOS DENSOS}
\end{center}

\vspace{0.5cm}

A presente invenção refere-se a método e sistema de manipulação
controlada do campo psiônico $\Psi$ mediante ressonância
eletromagnética na frequência do radical hidroxila ($\nuOH
\approx 1{,}666$ GHz, $\lambda \approx 18$ cm) acoplada a
regime criogênico ($T \leq \Tcrit = h\nuOH/k_B \approx
80$ mK). A frequência OH é identificada como frequência de
ressonância do campo $\Psi$ (``água escura''), com base na
correspondência entre OH (precursor molecular da água) e o
campo psiônico (substrato não-colapsado da matéria ordinária).
O método opera em três etapas: (1)~irradiação do meio-alvo em
$\nuOH$ com cavidade ressonante de alto $Q$;
(2)~resfriamento criogênico a $\leq 80$ mK por refrigeração de
diluição; (3)~nucleação de domínios de fase densa (``gotículas
psiônicas'', $\sim 100$ nm, $\rho \approx 2\rho_0$) onde o
acoplamento gravitacional-eletromagnético ($\alphaii = 0{,}012$,
Constante de Miguel) se torna operante. O sistema compreende
fonte de micro-ondas estabilizada, cavidade ressonante
($Q > 10^6$), criostato de diluição ($\leq 10$ mK), alvo
intercambiável e detectores multimodais. Aplicável a:
intensificação de reações nucleares D-D por redução de barreira
Coulombiana; indução de supercondutividade por coerência
psiônica; detecção de matéria escura por deslocamento de linhas
OH ($\Delta\nu/\nu \propto \alphaii \cdot \rhoPsi$); e
engenharia de materiais com propriedades anômalas.

\vspace{0.5cm}
\textbf{Figura principal: Figura 1} (Aparato completo)

\end{document}
